\section{Background}\label{sec:rs-background}

The problem of accountable recursive agent spawning sits at the intersection of four research traditions --- agent lifecycle management, accountability and provenance in distributed AI systems, delegation chain architecture, and hierarchical task decomposition --- and depends architecturally on the \bxthree{} three-layer model as its enforcement substrate.  Each tradition has produced significant contributions, but each also contains a specific structural gap that Recursive Spawning is designed to fill.  This section surveys all four areas, establishes exactly where those gaps lie, and situates the \bxthree{} Framework as the minimal sufficient substrate for immutable parent chains.

% ---------------------------------------------------------------
\subsection{Agent Lifecycle Management in Multi-Agent Systems}\label{sec:lifecycle}

Modern AI agent frameworks have moved away from treating agent creation as a singular one-time initialization event.  Instead, agent creation, assignment, delegation, revocation, and termination are understood as managed transitions within a structured lifecycle \cite{intelligent_ai_delegation}.  The Google DeepMind ``Intelligent AI Delegation'' framework identifies five core operations: \texttt{decompose} (breaking a goal into sub-goals), \texttt{delegate} (assigning a sub-goal to a sub-agent), \texttt{transfer} (reassigning authority across agents), \texttt{revoke} (withdrawing a delegation), and \texttt{audit} (reconstructing the delegation history for review).  These operations are treated not as independent primitives but as a composed system in which each has preconditions and postconditions that affect the others.

The structurally significant observation in this framework is that delegation is not merely one-hop assignment but can be \emph{recursive}.  The authors state explicitly that ``delegation can be recursive where an agent is assigned a task to identify and assign sub-tasks to others, effectively delegating the act of delegation itself'' \cite[intelligent_ai_delegation, \S4.2]{arxiv2602.11865}.  Furthermore, this decomposition may occur not only before task assignment but also dynamically during execution, meaning the depth and shape of the delegation tree are not fully specified at root-task initialization \cite[intelligent_ai_delegation, \S4.1]{arxiv2602.11865}.

This recursion pattern is the dominant operational mode in any non-trivial task environment.  A complex objective such as ``prepare the Q3 strategic acquisition report for the board'' decomposes into market analysis, financial modeling, legal review, and competitive intelligence --- each of which further decomposes into data collection, extraction, synthesis, and formatting tasks.  No single agent possesses the contextual capacity, tool access, and domain expertise to complete all of these without subdivision.  Recursive delegation is therefore not optional; it is the mechanism by which agent systems achieve practical coverage of complex goals.

The critical gap in current lifecycle management approaches is the treatment of the parent reference.  In most production multi-agent frameworks --- including CrewAI's hierarchical process mode \cite{zylos2026}, LangGraph's stateful graph runtimes \cite{heavybit2025multiagent}, and comparable systems --- the parent-child relationship is maintained as mutable runtime state.  In CrewAI's hierarchical process, the manager agent holds references to worker agents, but those worker references are updated dynamically as the orchestration progresses.  In LangGraph and similar graph-based runtimes, node edges are part of the graph state that can be rewritten during execution, meaning the delegation chain is a living data structure rather than a fixed record.

This mutability creates two distinct accountability problems.  First, the provenance trail is not immediately available; it is reconstructable only by replaying the entire graph state history, requiring both computational resources and complete runtime logs that most production systems do not retain at sufficient fidelity.  Second, any intermediate actor in the delegation chain can redirect the chain by overwriting the parent pointer, severing the connection between an action and its originating authority.  The DeepMind framework itself acknowledges that ``in long chains (A $\rightarrow$ B $\rightarrow$ C), accountability is transitive and provenance must be immutable'' \cite{intelligent_ai_delegation}.  Yet the framework's proposed solution --- registries listing agent skills and completion rates --- provides a nominal accountability record, not a structural one.  A registry can identify who acted after the fact; it cannot prevent a running agent from rewriting its parent reference to evade attribution in the first place.  Recursive Spawning addresses this gap at its root: the parent reference is made architecturally immutable, not merely recorded.

% ---------------------------------------------------------------
\subsection{Accountability and Provenance in Distributed AI Architectures}\label{sec:provenance}

The challenge of tracing actions to their originating authority in distributed AI systems has been treated in the literature primarily as a provenance problem.  Provenance, in the distributed systems sense, refers to the complete record of the causal chain connecting an output to its input sources and the agents that transformed those sources along the way \cite{on_chain_provenance_ai}.  The goal is to answer: \emph{who is responsible for this output, and through what chain of delegations did it come to pass?}  In conventional distributed systems, provenance is achieved through append-only logging: each event is hash-linked to its predecessor, making retroactive modification computationally infeasible without detection.

This approach has been extended to AI audit contexts in several important recent works.  ChainScore Labs argues that only cryptographically verifiable, on-chain provenance for training data, model weights, and inference outputs can create the immutable audit trail required for safety, compliance, and trust in autonomous AI systems \cite{on_chain_provenance_ai}.  Running model inference on-chain or proving it with zero-knowledge technology (ZKML) creates an immutable, time-stamped record for every AI decision, transforming AI audits from ``opaque promises into verifiable, time-stamped records of model lineage and data origin'' \cite{on_chain_provenance_ai}.  FlexBlok similarly proposes that blockchain-based distributed ledgers can establish independently verifiable, tamper-evident records for AI governance --- recording model lifecycle events, logging AI decisions, and managing digital identities without requiring trust in any single operator \cite{blockchain_ai_governance}.

The Human Delegation Provenance (HDP) protocol, introduced in 2025--2026, is the work most directly relevant to Recursive Spawning's contribution.  HDP proposes a lightweight token-based scheme that captures human authorization decisions and propagates them through the delegation chain as first-class cryptographic artifacts \cite{hdp2025arxiv}.  The protocol's core insight is that human authorization is not a single permission granted at chain initialization but a per-decision artifact that must travel with each delegation event.  Every HDP token encodes the scope, duration, and constraints of the human's delegation decision, and cryptographic verification is performed at each chain traversal.

However, HDP has a complementary blind spot that Recursive Spawning directly addresses.  HDP does not structurally enforce the integrity of the delegation chain itself --- it assumes that the chain through which its tokens travel is trustworthy.  In a mutable-chain environment, a compromised or misbehaving agent can present a different parent identity than the one under which it was actually created, inserting a false ancestor into the chain before presenting it to HDP token verification.  The HDP token will verify correctly against the false chain, because the token's own cryptographic integrity is independent of whether the chain topology is authentic.  The combination of HDP-style token provenance with a structurally immutable parent chain closes both gaps simultaneously: HDP provides per-decision human authorization, and Recursive Spawning provides the structurally immutable chain substrate that HDP's verification assumes but does not itself enforce.

The forensic ledger maintained by the \bxthree{} \fact{} layer achieves the provenance properties that blockchain proponents seek for AI systems, but without the operational overhead of a public distributed ledger.  ChainScore argues that ``regulators can audit the immutable record without trusting any single party's logs; incidents are traceable'' \cite{on_chain_provenance_ai}.  The \bxthree{} forensic ledger has this same property: it is append-only within the framework's own execution environment, every entry is hash-chained to its predecessor, and any external auditor can recompute the hash chain to confirm continuity without relying on the system's own operators or internal logs.

% ---------------------------------------------------------------
\subsection{Fixed vs. Mutable Delegation Chains}\label{sec:fixed-v-mutable}

The architectural choice between fixed and mutable delegation chains has consequences that ramify across every accountability property of a multi-agent system.  A mutable delegation chain allows parent references to be updated during execution, providing genuine operational advantages: agents can be dynamically reassigned in response to workload shifts, chains can be rebalanced to respond to performance bottlenecks, and failed agents can be replaced by transparently redirecting their children to new parents.  These conveniences are the reason that essentially all production multi-agent systems today implement mutable chains.

The accountability cost of this convenience is severe.  When a parent reference can change after the fact, the chain of responsibility becomes history-dependent in a way that defeats audit.  An action taken at time $t$ may be attributable to one parent, but if the reference is later updated to point to a different agent, the same action becomes attributable to a different parent retroactively.  Reconstructing the authoritative attribution chain requires capturing the complete mutable state at every point in time --- a requirement that most systems fail to meet.  In practice, mutable-chain systems either do not capture this history or capture it at insufficient fidelity to support definitive attribution.

A fixed delegation chain establishes parent references at node creation and preserves them unchanged for the node's lifetime.  This rigidity forgoes the operational flexibilities described above, but it produces a chain that is self-certifying: any participant can reconstruct the authoritative chain at any time by walking parent pointers from any node to the root.  No runtime state capture is required, because the chain structure itself is the record.  This property --- that the chain of custody is inherent in the structure rather than stored separately --- is what makes fixed chains suitable for high-stakes accountability environments where the record must be trustworthy independent of the system's own operators.

The practical cost of fixed chains is real but manageable.  Load rebalancing in a fixed-chain system requires spawning new nodes rather than redirecting existing ones; a failed agent cannot be transparently replaced by reassigning its children.  When an agent fails in a Recursive Spawning chain, its children remain attached to their original parent, and the failure is surfaced through the escalation path rather than concealed by a transparent reassignment that would break the provenance chain.  In high-stakes environments --- legal reasoning, medical diagnosis, financial compliance, critical infrastructure orchestration --- temporary throughput reduction during a failure incident is far less costly than accountability loss.  The correct trade-off in these domains is to preserve the chain and surface the failure to human judgment, which is exactly what Recursive Spawning does.

Existing approaches to chain integrity management are uniformly policy-based.  Access-control lists, role-based security models, and audit logs are all enforced by the runtime's execution logic \cite{nist2023}.  The fundamental vulnerability of policy-based chain management is that it trusts the agent to enforce constraints against its own operational interests.  An agent that is close to completing a valuable task has a strong operational incentive to suppress escalation, exceed its token budget, or silently absorb a failure rather than surface it to the human anchor and risk termination or reassignment.  Policy-based constraints can be bypassed by the agent that is subject to them.  Recursive Spawning addresses this by making the parent pointer a structural property enforced by the \fact{} layer, not a policy parameter that the agent can rewrite.

% ---------------------------------------------------------------
\subsection{Hierarchical Task Decomposition in AI Systems}\label{sec:task-decomposition}

Hierarchical task decomposition is the technique by which a complex goal is recursively subdivided into sub-goals assignable to specialized agents, enabling multi-agent systems to achieve coverage of tasks that no single agent could complete.  This technique appears across multiple research traditions.  The Recursive Delegation Engine (RDE) formalizes it as a tree-growing process in which each level of decomposition produces a more specialized, narrower-purposed sub-delegation \cite{rde2025}.  AgentOrchestra provides a systematic framework for hierarchical multi-agent task solving in which a manager agent decomposes goals and assigns subtasks to specialists, iterating until quality criteria are satisfied \cite{agentorchestra2025}.  IBM's analysis observes that ``relying on a single AI agent for a complex, multi-step task means the agent will typically use an expensive, large frontier AI model for every step, even when not every step requires this level of sophistication; agent hierarchies create the flexibility to fit the right AI model to each step'' \cite{ibm2025hierarchical}.  This cost-efficiency argument alone has driven widespread adoption of hierarchical decomposition in production agent systems.

Galileo AI's survey of multi-agent coordination failures identifies three emergent risks that appear specifically when delegation chains grow deep without constraint \cite{multi_agent_coordination_strategies}.  First, goal coherence degrades as context diverges from the original intent: each level of delegation introduces a new agent with its own context interpretation, and the divergence compounds with depth.  Second, attribution ambiguity makes it impossible to trace outputs to specific decomposition decisions: when a deep chain produces a poor output, it is unclear whether the fault lies in the root decomposition, an intermediate delegation decision, or the terminal execution.  Third, system paralysis occurs when the coordination overhead of managing a deeply nested chain exceeds the system's operational capacity.  These risks are not design flaws in any particular framework; they are emergent properties of recursive decomposition without structural depth constraints.

Existing approaches to depth management are uniformly policy-based.  Timeouts, token budgets, maximum chain length parameters, and human review triggers are all configured as operational parameters in the agent runtime and enforced by the runtime's execution logic \cite{multi_agent_coordination_strategies}.  The fundamental vulnerability of policy-based depth management is that it trusts the agent to enforce constraints against its own operational interests.  An agent near completion of a valuable task has a strong incentive to suppress escalation, exceed its token budget, or silently absorb a failure rather than surface it and risk termination.  Recursive Spawning addresses the depth problem structurally: the depth bound $D_{\max}$ is not a runtime parameter that the agent can modify but a constraint enforced by the \bxthree{} Bounds Engine layer, which operates in a separate execution context from the agents it constrains, so that no agent can modify the enforcement logic even indirectly.

% ---------------------------------------------------------------
\subsection{The BX3 Framework as Substrate for Immutable Parent Chains}\label{sec:bx3-substrate}

The \bxthree{} Framework is an accountability architecture for multi-agent AI systems that organizes every node into three immutable functional layers: the \purpose{} layer, the \bounds{} layer, and the \fact{} layer \cite{bx3framework}.  The \purpose{} layer defines authorization policy, establishes the human anchor $h(n)$ for each node, and receives all escalated exceptions.  The \bounds{} layer evaluates operating range constraints including the depth bound $D_{\max}$, the spawn necessity criterion, and resource allocation limits.  The \fact{} layer maintains the append-only forensic ledger, enforces the immutable parent pointer via pre-commit hooks, and executes the structural invariants that make all other constraints enforceable.

The framework's five pillars form an interdependent accountability system whose correctness properties require the simultaneous operation of all five \cite{bx3framework}.  Pillar 1 (Loop Isolation) ensures that no two nodes share mutable execution state without mediated communication through the \fact{} layer, preventing state corruption and ensuring that the provenance of every decision is traceable.  Pillar 2 (Recursive Spawning with Immutable Parent Pointers) establishes the parent reference as a structurally immutable property established at node creation.  Pillar 3 (Spatial Firewall) prevents cross-domain state contamination between nodes operating in distinct domains.  Pillar 4 (Sandbox Gate) enforces the operating range constraints of the \bounds{} layer as a structural gate rather than a policy parameter.  Pillar 5 (Bailout Protocol) closes the accountability loop by compulsorily escalating any unresolved exception to the human anchor $h(n)$ through the immutable parent chain, bypassing all intermediate machine actors \cite{bx3framework}.

The immutable parent pointer is Pillar 2 of the \bxthree{} Framework.  It is not a policy convention, not an operational guideline, and not a trust convention.  It is a structural requirement enforced by the \fact{} layer at node creation via a pre-commit hook that seals the parent reference before the node's execution context is initialized.  After this point, the reference is cryptographically bound to the node's identity and cannot be modified by any operation, including operations initiated by the node itself, by administrative actors, or by any other framework component.  This structural immutability is the propagation substrate for the Bailout Protocol: exception states have a structurally fixed escalation path through the immutable parent chain that no machine actor can redirect.

The three-layer architecture is the minimal structure that provides Recursive Spawning's correctness guarantees.  Removing the \purpose{} layer eliminates the exclusive terminus of escalation, allowing machine actors to absorb exceptions and issue terminal directives, causing the chain to terminate in autonomous resolution rather than human judgment.  Removing the \bounds{} layer converts the depth bound from a structural enforcement into a policy parameter that operational urgency can exceed.  Removing the \fact{} layer makes the immutable parent pointer unenforceable; the parent reference becomes mutable and the forensic ledger cannot be independently trusted.  This three-layer minimality is formally established in the \bxthree{} Framework's integration specification for Recursive Spawning, which proves that no proper subset of layers achieves the same correctness guarantees \cite{bx3framework}.

The \bxthree{} forensic ledger provides the append-only, hash-chained record of all spawn events, escalation events, and resolution directives that is the prerequisite for independently verifiable provenance.  Maintained by the \fact{} layer and independently verifiable through hash-chain recomputation, the ledger requires no trusted operator and no public blockchain infrastructure.  Its properties match what blockchain-based AI provenance systems seek to provide \cite{on_chain_provenance_ai}, but achieved within the framework's own execution environment with lower operational overhead and no dependency on external distributed consensus.  This makes the ledger suitable for regulatory audit, legal discovery, and adversarial review in high-stakes deployment environments.
